
Pátý případ užití je specifikován v tabulce č. \ref{tab:usecase:UC05}.

\begin{UseCaseTable}
  {Evidence budovy}
  {UC05}
  {F05}
  {Zaregistrování budovy do evidence systému}
  {Provozovatel monitorovaného prostoru}
  {}
  {Uživatel je přihlášen}
  {Budova je v systému evidována}

\UseCaseScenario{Základní scénář}
\UseCaseStep{uživatel}{V administraci zvolí akci pro založení budovy.}
\UseCaseStep{systém}{Vrátí formulář pro evidenci budovy se základními a doplňujícími údaji.}
\UseCaseStep{uživatel}{Vyplní povinné základní údaje (název, adresa, typ), volitelně doplňující údaje (souřadnice, rok výstavby, rok poslední rekonstrukce, úroveň veřejného zobrazení polohy) a formulář odešle.}
\UseCaseStep{systém}{Ověří, že povinné údaje jsou vyplněny a všechny zadané údaje jsou v platném tvaru.}
\UseCaseStep{systém}{Uloží budovu do evidence a zobrazí potvrzení.}

\UseCaseScenario{Alternativní scénář -- neplatné údaje}
\UseCaseStepCustom{A1}{uživatel, systém}{shodné s kroky 1--3 základního scénáře.}
\UseCaseStepCustom{A2}{systém}{Zjistí chybu ve vyplnění formuláře (chybí povinné pole, neplatný formát hodnoty) a vyzve uživatele k její opravě.}
\UseCaseStepCustom{A3}{uživatel}{Opraví zadané údaje.}
\UseCaseStepCustom{A4}{systém}{Pokračuje krokem 4 základního scénáře.}

\end{UseCaseTable}
